% ---------------------------------------------------------------
% TEMPLATE TCC1 - SENAC SANTO AMARO
% ---------------------------------------------------------------
% Trabalho de Conclusão de Curso 1
% Elaborado para apresentação da proposta inicial do TCC
% Estrutura: 4 capítulos (Introdução, Revisão Bibliográfica,
%            Desenvolvimento, Referências Bibliográficas)
% ---------------------------------------------------------------

\documentclass[
    12pt,                % tamanho da fonte
    openright,           % capítulos começam em página ímpar
    oneside,             % para impressão em apenas um lado do papel
    a4paper,             % tamanho do papel
    english,             % idioma adicional
    brazil               % idioma principal
]{abntex2}

% ---------------------------------------------------------------
% PACOTES
% ---------------------------------------------------------------
\usepackage{lmodern}                % Fonte Latin Modern
\usepackage[T1]{fontenc}            % Seleção de códigos de fonte
\usepackage[utf8]{inputenc}         % Codificação do documento
\usepackage{lastpage}               % Usado pela Ficha catalográfica
\usepackage{indentfirst}            % Indenta o primeiro parágrafo
\usepackage{color}                  % Controle de cores
\usepackage{graphicx}               % Inclusão de gráficos
\usepackage{microtype}              % Melhorias de justificação
\usepackage{float}                  % Para posicionamento de figuras
\usepackage{booktabs}               % Tabelas profissionais
\usepackage{multirow}               % Células mescladas em tabelas
\usepackage{longtable}              % Tabelas longas
\usepackage{amsmath,amssymb}        % Símbolos matemáticos
\usepackage{pgfgantt}               % Diagramas de Gantt
\usepackage[brazilian,hyperpageref]{backref}  % Páginas com citações
\usepackage[alf]{abntex2cite}       % Citações ABNT

% ---------------------------------------------------------------
% CONFIGURAÇÕES DO PDF
% ---------------------------------------------------------------
\makeatletter
\hypersetup{
    pdftitle={\@title},
    pdfauthor={\@author},
    pdfsubject={TCC1 - SENAC Santo Amaro},
    pdfcreator={LaTeX with abnTeX2},
    pdfkeywords={tcc}{senac}{santo amaro}{trabalho de conclusão},
    colorlinks=true,        % false: links em caixas; true: links coloridos
    linkcolor=blue,         % cor dos links internos
    citecolor=blue,         % cor dos links para bibliografia
    filecolor=magenta,      % cor dos links para arquivos
    urlcolor=blue,
    bookmarksdepth=4
}
\makeatother

% ---------------------------------------------------------------
% CONFIGURAÇÕES DE APARÊNCIA DO PDF FINAL
% ---------------------------------------------------------------
\definecolor{blue}{RGB}{41,5,195}

% ---------------------------------------------------------------
% INFORMAÇÕES DO DOCUMENTO
% ---------------------------------------------------------------
\titulo{Otimização de overbooking em consultórios médicos utilizando teoria das
filas}
\autor{Caio Troiano Collino (1142491237), Thiago Pereira Lins da Silva (1142515180), Vinícius Vianna Egydio de Carvalho (1142446597)}
\local{São Paulo}
\data{2026}
\orientador{Prof. Afonso Lelis}
% \coorientador{Prof. Nome do Coorientador} % Descomente se houver coorientador

\instituicao{%
  SENAC -- Serviço Nacional de Aprendizagem Comercial
  \par
  Unidade Santo Amaro
  \par
  Curso de Ciência da Computação
}

\tipotrabalho{Trabalho de Conclusão de Curso I}

\preambulo{Proposta de Trabalho de Conclusão de Curso apresentado ao SENAC Santo Amaro como requisito parcial para aprovação na disciplina de TCC1 do curso de Ciência da Computação.}

% ---------------------------------------------------------------
% COMPILA O ÍNDICE
% ---------------------------------------------------------------
\makeindex

% ---------------------------------------------------------------
% INÍCIO DO DOCUMENTO
% ---------------------------------------------------------------
\begin{document}

% Seleciona o idioma do documento
\selectlanguage{brazil}

% Retira espaço extra obsoleto entre as frases
\frenchspacing

% ---------------------------------------------------------------
% ELEMENTOS PRÉ-TEXTUAIS
% ---------------------------------------------------------------
\pretextual

% ---------------------------------------------------------------
% CAPA
% ---------------------------------------------------------------
\imprimircapa

% ---------------------------------------------------------------
% FOLHA DE ROSTO
% ---------------------------------------------------------------
\imprimirfolhaderosto

% ---------------------------------------------------------------
% RESUMO
% ---------------------------------------------------------------
\begin{resumo}
Este documento apresenta a proposta inicial do Trabalho de Conclusão de Curso (TCC1) desenvolvido no SENAC Santo Amaro. O objetivo deste trabalho é reduzir a taxa de evasão e melhorar a experiência do usuário em clínicas médicas, ao mesmo tempo mantendo o lucro. A justificativa para este estudo baseia-se em [descrever brevemente a justificativa]. A revisão bibliográfica abrange [mencionar os principais tópicos]. Como solução proposta, pretende-se realizar a otimização do overbooking a partir do uso de simulação, modelos probabilísticos e teoria das filas. Este trabalho está organizado em três capítulos que apresentam a introdução ao tema, a revisão bibliográfica que fundamenta a pesquisa, o desenvolvimento com a solução proposta e o cronograma de trabalho, seguidos das referências bibliográficas.

\vspace{\onelineskip}

\noindent
\textbf{Palavras-chave}: palavra1. palavra2. palavra3. palavra4. palavra5.
\end{resumo}

% ---------------------------------------------------------------
% LISTA DE ILUSTRAÇÕES
% ---------------------------------------------------------------
% \pdfbookmark[0]{\listfigurename}{lof}
% \listoffigures*
% \cleardoublepage

% ---------------------------------------------------------------
% LISTA DE TABELAS
% ---------------------------------------------------------------
% \pdfbookmark[0]{\listtablename}{lot}
% \listoftables*
% \cleardoublepage

% ---------------------------------------------------------------
% SUMÁRIO
% ---------------------------------------------------------------
\pdfbookmark[0]{\contentsname}{toc}
\tableofcontents*
\cleardoublepage

% ---------------------------------------------------------------
% ELEMENTOS TEXTUAIS
% ---------------------------------------------------------------
\textual

% ===============================================================
% CAPÍTULO 1 - INTRODUÇÃO
% ===============================================================
\chapter{Introdução}
\label{sec:introducao}

Overbooking é uma prática muito utilizada no setor aéreo, onde uma companhia vende mais passagens do que a capacidade real de assentos em um voo. O principal objetivo dessa prática é tentar garantir que o avião parta lotado, apesar de um certo número de passageiros não comparecerem (no-shows).

Apesar de ser uma prática mais conhecida pelo seu uso no setor da aviação, a mesma é vastamente utilizada no ambiente de clínicas e hospitais com o mesmo objetivo. Entretanto, segundo Zeng et al.\ (2013), a utilização de overbooking em serviços de atenção primária pode aumentar a taxa de ausência de pacientes e, para valores críticos do número de pacientes sob responsabilidade do médico (panel size), pode impactar negativamente o desempenho financeiro da clínica ou hospital.

\section{Justificativa}

\subsection{Relevância acadêmica}
% Escreva aqui a importância teórica do tema

\subsection{Relevância prática}
A otimização do overbooking pode contribuir para o aumento do lucro e para a redução do número de pacientes ausentes, melhorando a eficiência operacional de clínicas e hospitais.

\subsection{Lacuna identificada}
Ainda há limitações na compreensão dos impactos do overbooking em ambientes de clínicas e hospitais, especialmente no que diz respeito ao equilíbrio entre eficiência, qualidade do atendimento e experiência do paciente.

\section{Objetivos}

\subsection{Objetivo geral}
Otimizar o uso de overbooking em consultas médicas a partir da utilização de simulação, modelos probabilísticos e teoria das filas, com o objetivo de reduzir a taxa de evasão dos pacientes e melhorar sua experiência, mantendo a viabilidade econômica do sistema.

\subsection{Objetivos específicos}
\begin{itemize}
    \item Coletar dados a partir de um dataset e desenvolver um modelo baseado em conceitos probabilísticos;
    \item Simular o modelo e avaliar seu desempenho e viabilidade por meio de métodos estocásticos;
    \item Testar o modelo em um ou mais cenários de uso real;
    \item Analisar o funcionamento do mecanismo de overbooking em sistemas de saúde.
\end{itemize}

% ---------------------------------------------------------------
\chapter{Contexto}
\label{sec:contexto}

\section{Relevância acadêmica}

% ---------------------------------------------------------------
% ===============================================================
% CAPÍTULO 2 - REVISÃO BIBLIOGRÁFICA
% ===============================================================
\chapter{Revisão bibliográfica}

\section{Metodologia da pesquisa bibliográfica}

A pesquisa bibliográfica foi conduzida a partir da consulta a diferentes bases de dados acadêmicas e científicas, com o objetivo de identificar trabalhos relevantes sobre o tema.

\begin{itemize}
    \item \textbf{Bases de dados consultadas:} Google Scholar, Core, Taylor \& Francis Online, ProQuest;

    \item \textbf{Palavras-chave utilizadas:} \textit{queuing theory, overbooking, clinic, hospital, health, optimization, simulation};

    \item \textbf{Critérios de inclusão e exclusão:} foram considerados trabalhos publicados entre 2021 e 2025, nos idiomas português e inglês, incluindo artigos científicos e estudos de caso;

    \item \textbf{Período de abrangência:} 2021--2025.
\end{itemize}

% ---------------------------------------------------------------
\section{Referencial teórico}
\label{sec:referencial}

\subsection{Overbooking}

O \textit{overbooking} é a venda de um número de produtos ou serviços em maior número que o estoque atual ou a capacidade de serviço, sendo aplicado principalmente nas áreas da saúde e da aviação. O \textit{overbooking} surgiu como uma maneira de compensar o lucro perdido por conta de clientes que não comparecem ou cancelam o serviço (Denominado como \textit{no-shows}), utilizando todo o estoque ou capacidade do serviço disponível para balancear a eficiência operacional e a qualidade do serviço. No contexto clínico, o principal \textit{trade-off} do overbooking é entre a eficiência operacional e a qualidade do serviço: segundo Zeng et  al. (2013), o maior uso da equipe médica derivado do \textit{overbooking} pode trazer benefícios como maior lucro e redução do atraso para conseguir uma consulta (\textit{appointment-delay}). Porém, essa prática aumenta o tempo de espera no consultório para ser atendido (\textit{office-delay}) e pode até reduzir o lucro da clínica para valores críticos de \textit{panel-size}.

\subsection{Métricas}

Para avaliar o modelo, serão utilizadas métricas baseadas na qualidade e eficiência de serviço, levando em consideração o \textit{overbooking} e a estabilidade da clínica. A partir disso, escolhemos trabalhar com o tempo médio de \textit{office-delay}, o tempo do \textit{appointment-delay} e a taxa de \textit{no-shows}. Também serão utilizadas outras métricas como a taxa de utilização de recursos para garantir estabilidade.

\subsection{Rede Neural Recorrente}

Uma Rede Neural Recorrente (Também conhecido pela sigla RNN) é uma rede neural profunda, com diversas camadas escondidas, treinada a partir de dados sequenciais ou séries temporais, com o objetivo de criar um modelo de aprendizado de máquina para realizar predições. O principal diferencial das RNNs está em sua capacidade de utilizar informações de entradas prévias para influenciar o valor atual de entradas e saídas. Para uma determinada entrada, a RNN mantém um estado oculto, responsável por acompanhar o contexto e agir como uma memória que armazena informações das entradas anteriores, processando a entrada atual junto com o o estado oculto da etapa anterior em cada etapa para que o RNN seja capaz de "lembrar" de dados anteriores, influenciando o valor da saída atual. No contexto clínico, utilizaremos a RNN para prever a probabilidade de \textit{no-shows} e ajustar o nível de \textit{overbooking} em função disso.

% ---------------------------------------------------------------
\section{Metodologia da pesquisa bibliográfica}
\label{sec:metodologia_pesquisa}

A pesquisa bibliográfica foi realizada nas bases Google Scholar, IEEE Xplore e ACM Digital Library, utilizando as palavras-chave ``machine learning'', ``classificação de texto'' e ``processamento de linguagem natural''. Foram selecionados artigos publicados entre 2021 e 2026, em português e inglês, que abordassem diretamente técnicas de classificação automática de documentos.

\begin{itemize}
    \item \textbf{Bases de dados consultadas}: Google Scholar, IEEE Xplore, ACM Digital Library, SciELO, Periódicos CAPES, Semantic Scholar, Core e Taylor \& Francis Online;
    \item \textbf{Palavras-chave utilizadas}: liste os termos de busca em português e inglês;
    \item \textbf{Critérios de inclusão e exclusão}: quais critérios foram usados para selecionar ou descartar trabalhos (período, idioma, relevância, tipo de publicação);
    \item \textbf{Período de abrangência}: intervalo de anos das publicações consultadas.
\end{itemize}

\section{Trabalhos relacionados}
\label{sec:trabalhos_relacionados}

Apresente até \textbf{3 trabalhos} que \textbf{compartilhem o mesmo foco de resolução de problema} que o seu. O critério principal de seleção é: os trabalhos escolhidos devem ter tentado resolver o \textbf{mesmo problema} (ou um problema muito próximo) ao que você está propondo resolver. Não basta que o trabalho aborde o mesmo tema ou utilize a mesma tecnologia --- ele deve ter enfrentado o mesmo desafio prático ou teórico.

\textbf{Por que esse critério é importante?} Ao selecionar trabalhos que atacam o mesmo problema, você demonstra que:
\begin{itemize}
    \item O problema é relevante e já foi reconhecido por outros pesquisadores;
    \item Existem soluções anteriores, mas com limitações que o seu trabalho pretende superar;
    \item Seu trabalho não é uma repetição, mas um avanço em relação ao que já foi feito.
\end{itemize}

Para cada trabalho relacionado, descreva:
\begin{enumerate}
    \item \textbf{Qual problema} o trabalho tentou resolver (e por que é o mesmo problema que o seu);
    \item A metodologia ou abordagem utilizada para resolvê-lo;
    \item Os principais resultados obtidos;
    \item As limitações identificadas na solução proposta;
    \item Como o \textbf{seu trabalho se diferencia} ou avança em relação a ele.
\end{enumerate}

\subsection{Trabalho 1: [Título do trabalho]}

\citeonline{Silva2022} abordou o problema de [descrever o problema --- deve ser o mesmo ou muito similar ao seu]. Para resolvê-lo, os autores propuseram [descrever a solução]. A metodologia utilizada foi [descrever]. Os resultados mostraram que [descrever]. No entanto, a solução apresenta limitações como [descrever as lacunas que permanecem]. O presente trabalho diferencia-se por [explicar como sua proposta avança na resolução do mesmo problema].

\subsection{Trabalho 2: [Título do trabalho]}

[Descrever o segundo trabalho seguindo a mesma estrutura acima, sempre evidenciando que o problema abordado é o mesmo ou muito próximo ao seu.]

\subsection{Trabalho 3: [Título do trabalho]}

[Descrever o terceiro trabalho seguindo a mesma estrutura acima, sempre evidenciando que o problema abordado é o mesmo ou muito próximo ao seu.]

\vspace{0.5cm}

Uma tabela comparativa ajuda a sintetizar as diferenças entre os trabalhos:

\begin{table}[htb]
    \centering
    \caption{Comparação entre trabalhos relacionados}
    \label{tab:trabalhos_relacionados}
    \begin{tabular}{l|c|c|c|c}
        \toprule
        \textbf{Critério} & \textbf{Trabalho 1} & \textbf{Trabalho 2} & \textbf{Trabalho 3} & \textbf{Este trabalho} \\
        \midrule
        Objetivo           & [resumo]   & [resumo]   & [resumo]   & [resumo] \\
        \midrule
        Tecnologia         & [tecn.]    & [tecn.]    & [tecn.]    & [tecn.] \\
        \midrule
        Validação          & [método]   & [método]   & [método]   & [método] \\
        \midrule
        Limitação          & [limit.]   & [limit.]   & [limit.]   & --- \\
        \bottomrule
    \end{tabular}
    \legend{Fonte: Elaborado pelo autor (2025)}
\end{table}

% ===============================================================
% CAPÍTULO 3 - DESENVOLVIMENTO
% ===============================================================
\chapter{Desenvolvimento}
\label{cap:desenvolvimento}

Neste capítulo, apresenta-se a solução proposta para o problema identificado, bem como o cronograma de trabalho previsto para o desenvolvimento completo do projeto no TCC2.

% ---------------------------------------------------------------
\section{Solução proposta}
\label{sec:solucao}

Descreva detalhadamente a solução que será desenvolvida para resolver o problema apresentado na introdução. A descrição deve incluir:

\begin{itemize}
    \item \textbf{Visão geral da solução}: o que será desenvolvido (sistema, aplicativo, modelo, ferramenta, etc.);
    \item \textbf{Arquitetura}: como os componentes da solução se organizam e interagem;
    \item \textbf{Tecnologias e ferramentas}: quais linguagens, frameworks, bibliotecas, bancos de dados e ambientes serão utilizados, e por que foram escolhidos;
    \item \textbf{Funcionalidades principais}: o que a solução fará para atender aos objetivos;
    \item \textbf{Metodologia de desenvolvimento}: qual abordagem será utilizada (ágil, cascata, prototipação, etc.);
    \item \textbf{Critérios de validação}: como será medido o sucesso da solução (métricas, testes, avaliação com usuários, etc.).
\end{itemize}

Utilize diagramas, fluxogramas ou mockups para ilustrar a solução. Exemplo de descrição de arquitetura:

\textit{A solução proposta consiste em uma aplicação web composta por três camadas: front-end desenvolvido em React.js, back-end em Node.js com Express, e banco de dados PostgreSQL. A comunicação entre front-end e back-end será feita por meio de uma API REST.}

% ---------------------------------------------------------------
\section{Cronograma de trabalho}
\label{sec:cronograma}

O cronograma a seguir apresenta, no formato de diagrama de Gantt, as atividades previstas para o desenvolvimento do TCC2, distribuídas ao longo das semanas do semestre.

\begin{figure}[htb]
    \centering
    \caption{Cronograma de atividades -- Diagrama de Gantt}
    \label{fig:gantt}
    \begin{ganttchart}[
        hgrid,
        vgrid,
        x unit=0.72cm,
        y unit title=0.6cm,
        y unit chart=0.7cm,
        title height=1,
        bar height=0.5,
        bar label font=\scriptsize,
        title label font=\small,
        milestone label font=\scriptsize\itshape,
        group label font=\scriptsize\bfseries,
        bar/.append style={fill=blue!40},
        group/.append style={draw=black, fill=blue!70},
        milestone/.append style={fill=red!70},
    ]{1}{16}
        % Título: semanas de 1 a 16
        \gantttitle{Semanas}{16} \\
        \gantttitlelist{1,...,16}{1} \\

        % Grupo 1: Revisão e Planejamento
        \ganttgroup{Fase 1: Planejamento}{1}{4} \\
        \ganttbar{Revisão bibliográfica}{1}{3} \\
        \ganttbar{Definição de requisitos}{2}{4} \\
        \ganttbar{Projeto da arquitetura}{3}{4} \\
        \ganttmilestone{Entrega do projeto}{4} \\

        % Grupo 2: Desenvolvimento
        \ganttgroup{Fase 2: Desenvolvimento}{5}{11} \\
        \ganttbar{Implementação do módulo 1}{5}{7} \\
        \ganttbar{Implementação do módulo 2}{7}{9} \\
        \ganttbar{Integração dos módulos}{9}{11} \\
        \ganttmilestone{Protótipo funcional}{11} \\

        % Grupo 3: Testes e Finalização
        \ganttgroup{Fase 3: Validação}{12}{16} \\
        \ganttbar{Testes e validação}{12}{14} \\
        \ganttbar{Ajustes e correções}{14}{15} \\
        \ganttbar{Redação final do TCC2}{13}{16} \\
        \ganttmilestone{Entrega final}{16}

        % Ligações entre atividades
        \ganttlink{elem2}{elem3}
        \ganttlink{elem3}{elem4}
        \ganttlink{elem4}{elem5}
        \ganttlink{elem7}{elem8}
        \ganttlink{elem8}{elem9}
        \ganttlink{elem9}{elem10}
        \ganttlink{elem12}{elem13}
    \end{ganttchart}
    \legend{Fonte: Elaborado pelo autor (2025)}
\end{figure}

\textbf{Instruções para personalização do cronograma}:
\begin{itemize}
    \item Ajuste o número de semanas conforme o calendário acadêmico do semestre;
    \item Modifique as atividades de acordo com as etapas do seu projeto;
    \item Os marcos (\textit{milestones}) indicam entregas importantes --- adapte-os às datas de avaliação;
    \item As setas entre atividades indicam dependências --- uma atividade só inicia após a conclusão da anterior;
    \item Utilize o diagrama como ferramenta de acompanhamento: atualize-o a cada reunião com o orientador.
\end{itemize}

% ---------------------------------------------------------------
% ELEMENTOS PÓS-TEXTUAIS
% ---------------------------------------------------------------
\postextual

% ===============================================================
% CAPÍTULO 4 - REFERÊNCIAS BIBLIOGRÁFICAS
% ===============================================================
% As referências bibliográficas são geradas automaticamente a partir
% do arquivo bibliografia.bib, seguindo o padrão ABNT (NBR 6023).
%
% COMO CITAR no texto:
%   \cite{chave}         -> (SOBRENOME, ano) — citação indireta
%   \citeonline{chave}   -> Sobrenome (ano) — citação direta no texto
%
% TIPOS DE REFERÊNCIA no arquivo .bib:
%   @book        -> Livros
%   @article     -> Artigos de periódicos
%   @inproceedings -> Artigos em anais de eventos
%   @mastersthesis -> Dissertações de mestrado
%   @phdthesis   -> Teses de doutorado
%   @misc        -> Sites, vídeos e outros materiais
%   @manual      -> Documentação técnica
%
% Veja o arquivo bibliografia.bib para exemplos de cada tipo.
% ---------------------------------------------------------------
\bibliography{bibliografia}

% ---------------------------------------------------------------
% APÊNDICES (se necessário)
% ---------------------------------------------------------------
% \begin{apendicesenv}
% \partapendices
%
% \chapter{Nome do Apêndice}
% \label{ap:apendice1}
%
% Conteúdo do apêndice (material elaborado pelo próprio autor).
%
% \end{apendicesenv}

% ---------------------------------------------------------------
% ANEXOS (se necessário)
% ---------------------------------------------------------------
% \begin{anexosenv}
% \partanexos
%
% \chapter{Nome do Anexo}
% \label{an:anexo1}
%
% Conteúdo do anexo (material de terceiros).
%
% \end{anexosenv}

\end{document}
